\subsection{Circular Orbits and Kepler's Third Law}

For a satellite in a circular orbit of radius $r$ around a body with
gravitational parameter $\mu = GM$, Newton's law and the centripetal
acceleration condition give the orbital speed:
\begin{equation}
  v_c = \sqrt{\frac{\mu}{r}}.
  \label{eq:vc}
\end{equation}
The orbital period follows from circumference divided by speed:
\begin{equation}
  T = \frac{2\pi r}{v_c} = 2\pi\sqrt{\frac{r^3}{\mu}}.
  \label{eq:T}
\end{equation}
Equation~\eqref{eq:T} encapsulates the key insight for orbit phasing:
\textbf{lower orbit $\Rightarrow$ faster speed, shorter period;
higher orbit $\Rightarrow$ slower speed, longer period.}

\subsection{Argument of Latitude}

The \emph{argument of latitude} $u$ is the angle from the ascending node
$\Omega$ to the satellite's current position, measured along the orbital
plane. It uniquely describes the satellite's location on a circular orbit
and is the quantity we seek to optimise.

\subsection{Orbit Phasing}

A two-burn \emph{phasing manoeuvre} shifts a satellite's argument of
latitude by $\du$ without permanently altering the orbital radius.

\paragraph{Advancing ($\du > 0$).}
\begin{enumerate}
  \item Fire \textbf{retrograde} (opposite the velocity). The satellite
        slows and drops to a lower phasing orbit.
  \item The shorter period causes the satellite to lap the reference
        position in one revolution.
  \item Fire \textbf{prograde} at the same orbital location to restore
        the original circular orbit.
\end{enumerate}

\paragraph{Falling back ($\du < 0$).}
\begin{enumerate}
  \item Fire \textbf{prograde}. The satellite accelerates and climbs to
        a higher phasing orbit.
  \item The longer period means it completes one revolution after the
        reference position.
  \item Fire \textbf{retrograde} to restore the orbit.
\end{enumerate}

Both cases apply two burns of equal magnitude at the same orbital
location (the intersection of the circular orbit and the phasing ellipse).

\subsection{Delta-$v$ Cost}

To advance by angle $\du$ in exactly one revolution, the phasing orbit
period must satisfy $T_\text{ph}/T_c = 1 - \du/(2\pi)$.
Applying Kepler's third law~\eqref{eq:T} gives the phasing orbit
semi-major axis:
\begin{equation}
  a_\text{ph} = \frac{r}{\bigl(1 + \du/(2\pi)\bigr)^{2/3}}.
  \label{eq:aph}
\end{equation}
The speed at the manoeuvre point on the phasing orbit
(vis-viva equation):
\begin{equation}
  v_\text{ph} = \sqrt{\mu\!\left(\frac{2}{r} - \frac{1}{a_\text{ph}}\right)}.
  \label{eq:vph}
\end{equation}
Since both burns have equal magnitude, the total $\dv$ per satellite is:
\begin{equation}
  \dv_i = 2\,|v_\text{ph}(\du_i) - v_c|.
  \label{eq:dv}
\end{equation}
For $\du > 0$, equation~\eqref{eq:aph} gives $a_\text{ph} < r$ (first
burn retrograde); for $\du < 0$, $a_\text{ph} > r$ (first burn prograde).
